\documentclass[11pt,reqno]{amsart}
\usepackage[T1]{fontenc}
\usepackage{lmodern}
\usepackage{microtype}
\usepackage{mathtools,amssymb,amsfonts}
\usepackage{mathrsfs}
\usepackage[margin=1in]{geometry}
\usepackage{enumitem}
\usepackage{xurl}
\usepackage[hidelinks,pdfencoding=auto,psdextra]{hyperref}
\usepackage{aliascnt}
\usepackage[nameinlink,noabbrev]{cleveref}
\hypersetup{
  pdftitle={Odd polynomial counterexamples for the Ad root system in every dimension d >= 4},
  pdfauthor={AI-generated verification draft},
  pdfsubject={Proposed proof for independent human verification}
}
\allowdisplaybreaks[2]
\setlength{\emergencystretch}{2em}
\setlist[enumerate]{leftmargin=2.2em,itemsep=0.3em,topsep=0.4em}
\newtheorem{theorem}{Theorem}[section]
\newaliascnt{proposition}{theorem}
\newtheorem{proposition}[proposition]{Proposition}
\aliascntresetthe{proposition}
\crefname{proposition}{proposition}{propositions}
\Crefname{proposition}{Proposition}{Propositions}
\newaliascnt{lemma}{theorem}
\newtheorem{lemma}[lemma]{Lemma}
\aliascntresetthe{lemma}
\crefname{lemma}{lemma}{lemmas}
\Crefname{lemma}{Lemma}{Lemmas}
\newaliascnt{corollary}{theorem}
\newtheorem{corollary}[corollary]{Corollary}
\aliascntresetthe{corollary}
\crefname{corollary}{corollary}{corollaries}
\Crefname{corollary}{Corollary}{Corollaries}
\theoremstyle{definition}
\newaliascnt{definition}{theorem}
\newtheorem{definition}[definition]{Definition}
\aliascntresetthe{definition}
\crefname{definition}{definition}{definitions}
\Crefname{definition}{Definition}{Definitions}
\theoremstyle{remark}
\newaliascnt{remark}{theorem}
\newtheorem{remark}[remark]{Remark}
\aliascntresetthe{remark}
\crefname{remark}{remark}{remarks}
\Crefname{remark}{Remark}{Remarks}
\numberwithin{equation}{section}
\newcommand{\R}{\mathbb R}
\newcommand{\C}{\mathbb C}
\newcommand{\one}{\mathbf 1}
\newcommand{\eps}{\varepsilon}
\newcommand{\Orth}{\mathcal O_d}
\newcommand{\Pol}{\mathcal P_d}
\newcommand{\Exc}{\mathcal E}
\newcommand{\ip}[2]{\langle #1,#2\rangle}
\newcommand{\norm}[1]{\lVert #1\rVert}
\newcommand{\abs}[1]{\lvert #1\rvert}
\DeclareMathOperator{\Sym}{Sym}
\DeclareMathOperator{\diag}{diag}
\DeclareMathOperator{\supp}{supp}
\DeclareMathOperator{\spec}{spec}
\DeclareMathOperator{\RePart}{Re}
\DeclareMathOperator{\ImPart}{Im}
\title[Odd polynomial counterexamples for the $A_d$ root system]
{Odd polynomial counterexamples for the $A_d$ root system\\
in every dimension $d\geq4$}
\author{}
\date{September 4, 2026}
\keywords{Root systems, universal covers, odd polynomials, circle symmetry,
semialgebraic dimension}

\begin{document}
\begin{abstract}
For each $d\geq4$, we give an existence proof of an odd real polynomial
$F$ of degree at most $3d(d+1)-3$ whose restriction to the unit sphere
agrees with no linear functional on any orthogonal image of the normalized
$A_d$ root system. The construction adds an arbitrarily small
circle-invariant polynomial perturbation to the cubic--quintic seed of a
previously proposed dimension-four and dimension-five construction. A rank-one
spectral argument bounds the dimension of every signed circle-orbit
stratum of root configurations. Invariant polynomial interpolation and an
incidence-dimension argument then exclude all configurations having no
root in either of the two nontrivial weight planes. The remaining
configurations are treated by a necessary normal form and an explicit,
nonzero quintic obstruction. The proof does not require a
computer-algebra certificate or a classification of all cubic solutions.
We also derive the corresponding counterexamples to the regular-simplex
universal-cover conjecture. The auxiliary perturbation is existential;
we do not prove that the unmodified quintic seed works in all dimensions.
\end{abstract}
\maketitle

\noindent\textbf{Verification status.}
This is an AI-generated manuscript. It presents a proposed proof
for independent human verification; it has not been independently verified
or peer reviewed. The algebraic audit supplied with the source checks
selected identities, not the validity of the entire proof.

\tableofcontents
\clearpage

\section{Statement, provenance, and proof strategy}\label{sec:introduction}

Set $N=d+1$, and give
\[
 V=\one^\perp\subset\R^N,\qquad \one=(1,\ldots,1)^T,
\]
its induced Euclidean metric. The normalized positive roots are
\begin{equation}\label{eq:roots}
 \alpha_{ij}=\frac{e_i-e_j}{\sqrt2},\qquad 0\leq i<j\leq d.
\end{equation}
Their number is $n=\binom N2=d(d+1)/2$. Normalization matters here:
\emph{all root-label evaluations are at unit vectors}. Auxiliary polynomial
and gradient evaluations at the simplex vertices $p_i$ below are not unit-vector
evaluations; in fact $\norm{p_i}^2=d/(d+1)$.
Let $\Orth$ be the space of orthogonal isomorphisms $V\to\R^d$.
After a fixed orthogonal identification of these spaces, $\Orth$ is
$O(d)$. Superscript $T$ denotes the Euclidean adjoint. We write
$\Sym(V)$ for the real vector space of self-adjoint endomorphisms of $V$,
and $cc^T$ for the operator $w\mapsto\ip c w\,c$.

\begin{theorem}\label{thm:main}
For every integer $d\geq4$, there exists an odd real polynomial
$F:\R^d\to\R$ of degree at most
\begin{equation}\label{eq:degree-bound}
 D=6n-3=3d(d+1)-3
\end{equation}
such that no $A\in\Orth$ and $b\in\R^d$ satisfy
\begin{equation}\label{eq:agreement}
 F(A\alpha_{ij})=\ip{b}{A\alpha_{ij}}
 \qquad(0\leq i<j\leq d).
\end{equation}
In particular, $F|_{S^{d-1}}$ is a smooth odd counterexample even when
only orientation-preserving placements are allowed.
\end{theorem}

There is a more precise perturbative conclusion. Fix the orthogonal
splitting
\begin{equation}\label{eq:splitting}
 \R^d=E_1\oplus E_2\oplus E_0
      =\C\oplus\C\oplus\R^{d-4}
\end{equation}
and define
\begin{align}
 h_\theta(z_1,z_2,t)&=(e^{i\theta}z_1,e^{2i\theta}z_2,t),
       \label{eq:circle}\\
 Q(z_1,z_2,t)&=\RePart(z_1^2\overline z_2),\qquad
 H(z_1,z_2,t)=|z_1|^2\ImPart(z_1^2\overline z_2).
       \label{eq:seed}
\end{align}
Let $\Pol$ denote the real vector space of odd polynomials of degree
at most $D$ that are invariant under every $h_\theta$.

\begin{theorem}\label{thm:perturbative}
For each $d\geq4$, there is $\eps_0>0$ with the following property.
For every $0<|\eps|<\eps_0$ and every $\eta>0$, there exists
$R\in\Pol$, with $\norm R<\eta$ in any fixed norm on $\Pol$, such that
\[
 F=Q+\eps H+R
\]
has no agreement of the form \eqref{eq:agreement}.
\end{theorem}

For the geometric interpretation, define
\begin{equation}\label{eq:cover-polytope}
 C_d=\{y\in V:|\ip{y}{\alpha_{ij}}|\leq\tfrac12
                 \text{ for every }i<j\}.
\end{equation}
To specify the scale, put
$\Delta=\operatorname{conv}\{e_i-N^{-1}\one\}$ and use the polar in $V$,
$L^\circ=\{y\in V:\ip y x\leq1\text{ for every }x\in L\}$.
Since $\Delta-\Delta=\operatorname{conv}\{e_i-e_j:i\ne j\}$,
\[
 C_d=\frac1{\sqrt2}(\Delta-\Delta)^\circ.
\]
Thus the conjectured cover is the scaled \emph{polar of the difference
body}, not the difference body itself. Every defining hyperplane in
\eqref{eq:cover-polytope} is a facet hyperplane at distance $1/2$
from the origin: its tangency point with the radius-$1/2$ ball satisfies
all other defining inequalities strictly.

Here a \emph{universal cover for sets of diameter one} means a set
containing a congruent copy of every such set; equivalently, for each
such set $K$, some translated orthogonal image of the cover contains $K$.
Reflections as well as rotations and translations are allowed. In this
normalization, Makeev's regular-simplex conjecture asserts that $C_d$ is
a universal cover; see \cite[Conjecture 1]{Kuperberg1999}.

\begin{corollary}\label{cor:cover}
For every $d\geq4$, there is a compact convex body $K\subset\R^d$
of constant width $1$ that is not contained in $A C_d+b$ for any
$A\in\Orth$ and $b\in\R^d$. Thus $C_d$ is not a universal cover for
sets of diameter $1$.
\end{corollary}

\subsection{Prior work and attribution}
Kuperberg \cite[Section 1, Proposition 1]{Kuperberg1999} identifies
agreement of odd functions with linear functions on rotated facet
normals as the functional form of the constant-width circumscription
problem. His paper proves the positive result in dimension three and
explains why the corresponding homological obstruction vanishes in
higher dimensions. Its Section 5 also uses determination of quadratic
forms by root evaluations in the low-dimensional affine problem.
We give the required uniform linear-algebra statement below.

The specific polynomials $Q,H$ in \eqref{eq:seed}, the weight-$(1,2)$
circle action, the exact-label reduction, the cubic annihilator, and the
compactness perturbation strategy are inherited from the AI-generated
repository manuscript \cite{ChatGPT2026}, which reports counterexamples
in dimensions four and five. The VibeMathed entry \cite{VibeMathed2026}
is the route by which that manuscript was supplied. We reproduce the
identities needed here and do not use its dimension-specific algebraic
certificates or assume the correctness of those certificates.

Chen \cite[Section 2]{Chen2026} gives counterexamples for other facet
configurations. More particularly, his Appendix A, Proposition 13,
reproduces Makeev's finite-dimensional interpolation and dimension-count
argument above the threshold $n=d(d+1)/2$. The generic-avoidance argument
in this paper follows that general pattern. At the threshold,
unrestricted dimension counting alone gives no strict inequality.
Here we impose a circle symmetry, separate signed orbit types, and
exploit the one-dimensional symmetry in every incidence fiber.
The spectral fiber bound and the uniform treatment of exceptional
placements are the additional arguments proposed in this manuscript.
Chen's positive theorem at the maximal facet count concerns other
polytopes and a different conjecture of Makeev; it does not establish
the universal-cover property of $C_d$.

Our only non-elementary proof input is standard semialgebraic
triviality and dimension theory, recorded with references in
\cref{sec:dimension}. The complete-graph cut/cycle identity is also
related to the combinatorial Hodge decomposition in
\cite{JiangLimYaoYe2011}; the elementary part used here is proved directly.
No assertion of novelty is made for these background tools.

\subsection{Organization}
The proof has two independent parts. In \cref{sec:generic}, generic
polynomials in $\Pol$ exclude all placements for which no root lies in
$E_1\cup E_2$. In \cref{sec:cubic,sec:exceptional}, the seed
$Q+\eps H$ excludes the complementary compact set for sufficiently
small nonzero $\eps$. Openness on that compact set permits the two
parts to be combined in \cref{sec:completion}.
\Cref{sec:geometry} proves the geometric corollary, including the
support-function construction.

\section{Simplex frames, exact labels, and quadratic evaluations}
\label{sec:linear}

Put
\[
 q_i=e_i-\frac1N\one,\qquad
 \Pi=I_N-\frac1N\one\one^T.
\]
For $A\in\Orth$, write $p_i=Aq_i$.

\begin{lemma}\label{lem:frame}
The vectors $p_i$ satisfy
\begin{equation}\label{eq:frame}
 \sum_i p_i=0,\qquad
 \ip{p_i}{p_j}=\delta_{ij}-\frac1N,\qquad
 \sum_i p_ip_i^T=I_d.
\end{equation}
Consequently, the $d$ coordinate columns of the $N\times d$ matrix
with rows $p_i^T$ are orthonormal and perpendicular to $\one$.
\end{lemma}
\begin{proof}
The first two assertions follow by expanding $q_i=e_i-N^{-1}\one$.
Also $\sum_iq_iq_i^T=\Pi$, whose restriction to $V$ is the identity.
Apply the orthogonal map $A$ to obtain \eqref{eq:frame}. The last
assertion is the matrix form of the first and third identities.
\end{proof}

For an odd function $G$ on the unit sphere, define the skew-symmetric
label matrix $Y_G(A)$ by
\begin{equation}\label{eq:labels}
 Y_G(A)_{ij}=G\!\left(\frac{p_i-p_j}{\sqrt2}\right)
 \quad(i\ne j),\qquad Y_G(A)_{ii}=0.
\end{equation}
For polynomials the same notation refers to their restrictions to the
sphere. Define the residual
\begin{equation}\label{eq:residual}
 \mathcal R_G(A)=\Pi Y_G(A)\Pi.
\end{equation}

\begin{lemma}[Exact labels]\label{lem:exact}
For a skew-symmetric real $N\times N$ matrix $Y$, the following are
equivalent:
\begin{enumerate}[label=(\roman*)]
 \item $Y_{ij}=\ip{c}{\alpha_{ij}}$ for all $i<j$, for some $c\in V$;
 \item $Y_{ij}=\rho_i-\rho_j$ for real numbers $\rho_0,\ldots,\rho_d$;
 \item $\Pi Y\Pi=0$.
\end{enumerate}
Thus \eqref{eq:agreement} at a specified $A$ is equivalent to
$\mathcal R_F(A)=0$.
\end{lemma}
\begin{proof}
Condition (i) gives (ii) with $\rho_i=c_i/\sqrt2$. Conversely, subtract
the mean from the $\rho_i$ and take $c_i=\sqrt2\rho_i$. This gives
$c\in V$ and (i).

Under (ii), $Y=\rho\one^T-\one\rho^T$, so (iii) follows from
$\Pi\one=0$. For the converse put $J_0=N^{-1}\one\one^T$.
Skew-symmetry gives $J_0YJ_0=0$, and expanding (iii) yields
$Y=J_0Y+YJ_0$. Its entries are $\rho_i-\rho_j$, where
$\rho_i=N^{-1}\sum_jY_{ij}$.
Finally, in \eqref{eq:agreement} one has $c=A^Tb$, and every $c\in V$
arises from $b=Ac$.
\end{proof}

Matrices satisfying these conditions will be called \emph{exact}.
If $X$ is skew-symmetric and $X\one=0$, then
\begin{equation}\label{eq:cycle}
 \sum_{i<j}X_{ij}(\rho_i-\rho_j)
 =\sum_i\rho_i\sum_jX_{ij}=0.
\end{equation}
This is the cut/cycle orthogonality used later.

\begin{lemma}[Quadratic evaluation isomorphism]\label{lem:quadratic}
The linear map
\begin{equation}\label{eq:quadratic-map}
 \Sym(V)\longrightarrow\R^n,\qquad
 B\longmapsto\bigl(\ip{\alpha_{ij}}{B\alpha_{ij}}\bigr)_{i<j}
\end{equation}
is an isomorphism.
\end{lemma}
\begin{proof}
Extend $B$ to $\R^N$ by setting $B\one=0$. Form the symmetric matrix
$W$ with zero diagonal and entries
\[
 W_{ij}=\frac{B_{ii}+B_{jj}}2-B_{ij}
       =\ip{\alpha_{ij}}{B\alpha_{ij}}\quad(i\ne j).
\]
The same displayed expression is zero when $i=j$. Hence
$\Pi W\Pi=-B$, since $\Pi B\Pi=B$ and the terms involving $B_{ii}$
or $B_{jj}$ are killed by the projections. Thus \eqref{eq:quadratic-map}
is injective. Both spaces have dimension $d(d+1)/2=n$, proving
surjectivity as well.
\end{proof}

\section{Semialgebraic dimension facts}\label{sec:dimension}

We use real semialgebraic dimension, with $\dim\varnothing=-\infty$.
The needed foundational results are the Tarski--Seidenberg theorem,
semialgebraic dimension theory, and Hardt's triviality theorem;
see \cite[Sections 2.2, 2.8, and 9.3]{BCR1998} and \cite{Hardt1980}.
An accessible account of the stratification and trivialization
statements is \cite[Section 2.1]{HLTV2011}.

\begin{lemma}\label{lem:dimension}
Semialgebraic images are semialgebraic. Finite unions have dimension
equal to the maximum of the dimensions, and dimensions add under
Cartesian products. A semialgebraic subset of $\R^s$ of dimension
less than $s$ has empty interior.

Let $\pi:X\to Y$ be a continuous semialgebraic map. If every nonempty
fiber has dimension at most $k$, then
\begin{equation}\label{eq:fiber-upper}
 \dim X\leq\dim\pi(X)+k.
\end{equation}
If every nonempty fiber has dimension at least $k$, then
\begin{equation}\label{eq:fiber-lower}
 \dim X\geq\dim\pi(X)+k.
\end{equation}
In particular, when all nonempty fibers have dimension $k$, equality
holds.
\end{lemma}
\begin{proof}
The first assertions are the cited basic dimension results. By
semialgebraic triviality, $\pi(X)$ has a finite semialgebraic partition
$Y_a$ over which $\pi^{-1}(Y_a)$ is semialgebraically homeomorphic
to $Y_a\times F_a$, with $\pi$ corresponding to projection.
The two fiber inequalities follow from
\[
 \dim X=\max_a(\dim Y_a+\dim F_a),\qquad
 \dim\pi(X)=\max_a\dim Y_a.\qedhere
\]
\end{proof}

\section{Generic avoidance outside the exceptional set}\label{sec:generic}

\subsection{Orbit separation and interpolation}
The action \eqref{eq:circle} is semialgebraic: write
$\cos\theta=a$, $\sin\theta=b$ with $a^2+b^2=1$, and use
$\cos2\theta=a^2-b^2$, $\sin2\theta=2ab$.

\begin{definition}\label{def:exceptional}
A placement $A\in\Orth$ is \emph{exceptional} if $A\alpha_{ij}$
belongs to $E_1\cup E_2$ for at least one $i<j$. Write $\Exc$ for
the set of exceptional placements.
\end{definition}

The set $\Exc$ is a finite union of closed semialgebraic subsets of
the compact space $\Orth$, and hence is compact and semialgebraic.
It is nonempty, since an orthogonal placement can send any chosen root
to a unit vector in $E_1$. An odd circle-invariant function vanishes on
$S^{d-1}\cap(E_1\cup E_2)$, which explains why these placements must
be treated separately from the interpolation argument.

\begin{lemma}\label{lem:orbits}
For $w\in S^{d-1}$, the orbit of $w$ under \eqref{eq:circle} contains
$-w$ if and only if $w\in E_1\cup E_2$.
Moreover, the polynomial map
\begin{equation}\label{eq:orbit-map}
 \Psi(z_1,z_2,t)=
 \bigl(t,|z_1|^2,|z_2|^2,
        \RePart(z_1^2\overline z_2),
        \ImPart(z_1^2\overline z_2)\bigr)
\end{equation}
separates circle orbits in $\R^d$ and has degree at most three.
\end{lemma}
\begin{proof}
If $h_\theta w=-w$, the fixed component $t$ vanishes. If both complex
coordinates were nonzero, we would have simultaneously
$e^{i\theta}=-1$ and $e^{2i\theta}=-1$, which is impossible.
Conversely, $\theta=\pi$ reverses a vector of $E_1$ and
$\theta=\pi/2$ reverses a vector of $E_2$.

The entries of $\Psi$ are invariant. Suppose two points have the same
image. If their first complex coordinates are nonzero, their equal
moduli permit a unique phase rotation taking the first coordinate of
one point to that of the other. Equality of $z_1^2\overline z_2$
then gives equality of the second coordinates after that same
rotation. If both first coordinates vanish, equality of $|z_2|$
suffices, since multiplication by $e^{2i\theta}$ is transitive on
each circle. The fixed components are explicitly recorded. This
also covers the case of two zero complex coordinates.
\end{proof}

\begin{lemma}[Invariant interpolation]\label{lem:interpolation}
Let $w_1,\ldots,w_m\in S^{d-1}$, where $m\leq n$, be such that
the $2m$ circle orbits of $w_1,-w_1,\ldots,w_m,-w_m$ are distinct.
Then the evaluation map
\[
 \Pol\longrightarrow\R^m,\qquad
 F\longmapsto(F(w_1),\ldots,F(w_m))
\]
is surjective.
\end{lemma}
\begin{proof}
Let $\xi_1,\ldots,\xi_{2m}$ be the distinct points obtained by applying
$\Psi$ to the indicated signed points. For each $j$, define a polynomial
in the target coordinates $Y$ by
\[
 L_j(Y)=\prod_{k\ne j}
  \frac{\ip{Y-\xi_k}{\xi_j-\xi_k}}{\norm{\xi_j-\xi_k}^2}.
\]
It has degree at most $2m-1$ and satisfies $L_j(\xi_k)=\delta_{jk}$.
For arbitrary prescribed real values $a_1,\ldots,a_m$, a linear
combination $P$ of these polynomials takes the value $a_j$ at
$\Psi(w_j)$ and $-a_j$ at $\Psi(-w_j)$.
The polynomial $P\circ\Psi$ is circle-invariant and has degree
at most $3(2m-1)\leq D$. Its odd part
\[
 F(w)=\frac{P(\Psi(w))-P(\Psi(-w))}{2}
\]
is still invariant, belongs to $\Pol$, and has $F(w_j)=a_j$.
\end{proof}

\subsection{A rank-one spectral fiber bound}
For $\lambda=(\lambda_1,\lambda_2)\in(0,1)\times(4,5)$ set
\begin{equation}\label{eq:Dlambda}
 D_\lambda=\lambda_1I_{E_1}\oplus\lambda_2I_{E_2}
       \oplus\diag(\mu_1,\ldots,\mu_{d-4}),
 \qquad \mu_j=4j+4.
\end{equation}
The final block is absent when $d=4$. Distinct eigenvalues of
$D_\lambda$ are separated by more than three.

\begin{lemma}[Spectral fibers]\label{lem:spectral}
Every fiber of the semialgebraic map
\begin{equation}\label{eq:spectral-map}
 \Phi:\Orth\times(0,1)\times(4,5)\times\{c\in V:\norm c<1\}
 \longrightarrow\Sym(V),\qquad
 \Phi(A,\lambda_1,\lambda_2,c)=A^TD_\lambda A+cc^T
\end{equation}
has dimension at most two.
\end{lemma}
\begin{proof}
Fix a matrix $B$ in the image and write
\[
 S=A^TD_\lambda A,\qquad B=S+cc^T.
\]
In each two-dimensional eigenspace of $S$ for $\lambda_1$ or
$\lambda_2$, there is a nonzero vector perpendicular to $c$.
It remains an eigenvector of $B$ with the same eigenvalue.
Consequently $\lambda_1,\lambda_2\in\spec B$, leaving only finitely
many possibilities for the pair $\lambda$.

Fix such a pair. Write $s_1\leq\cdots\leq s_d$ for the eigenvalues
of $S$ and $\beta_1\leq\cdots\leq\beta_d$ for those of $B$,
with multiplicities. The inequalities
\[
 S\leq B\leq S+\norm c^2I
\]
in the order of symmetric matrices imply
\begin{equation}\label{eq:eigen-bounds}
 s_j\leq\beta_j\leq s_j+\norm c^2<s_j+1.
\end{equation}
For completeness, this follows from
\[
 \lambda_j(T)=\min_{\dim W=j}\ 
              \max_{\substack{x\in W\\\norm x=1}}\ip{x}{Tx}.
\]
The formula itself follows from an orthonormal eigenbasis: the span
of the first $j$ eigenvectors gives one inequality, and every
$j$-dimensional subspace meets the span of eigenvectors numbered
$j,\ldots,d$ nontrivially, giving the other.

The intervals $[s,s+1)$ attached to distinct eigenvalues $s$ of $S$
are disjoint. A simple eigenvalue of $S$ contributes only one ordered
eigenvalue of $B$ to its interval. A double eigenvalue contributes two,
and at least one equals the original eigenvalue. It follows that an
eigenvalue of $B$ has multiplicity at most two. More explicitly, a
repeated value must lie in an interval coming from a double eigenvalue
$s$ of $S$. That interval contains exactly two eigenvalues of $B$,
counting multiplicity, and at least one of them is $s$. If they coincide,
both are therefore $s$. Hence every double eigenvalue of $B$ is either
$\lambda_1$ or $\lambda_2$, and is also double in $S$.

Let $P_\beta$ be the orthogonal projection to the $\beta$-eigenspace
of $B$. The determinant identity $\det(I+uv^T)=1+v^Tu$, obtained
by multilinearity in the columns, gives the rational-function identity
\begin{equation}\label{eq:secular}
 \frac{\det(tI-S)}{\det(tI-B)}
 =1+c^T(tI-B)^{-1}c
 =1+\sum_{\beta\in\spec B}
       \frac{\norm{P_\beta c}^2}{t-\beta}.
\end{equation}
This identity is first obtained for $t\notin\spec B$ and then read as
an identity of rational functions. The left side is determined by $B$
and $\lambda$, because the spectrum of $S$ is prescribed by
$D_\lambda$. Its residues therefore determine every
$\norm{P_\beta c}^2$.
If $\beta$ is double in $B$, it is also double in $S$, so the left
side has no pole at $\beta$. Thus $P_\beta c=0$. On every remaining,
one-dimensional eigenspace, the component of $c$ is determined up to sign (with just one choice for a zero component). Thus there are
at most $2^d$ possible vectors $c$ for this fixed pair $\lambda$,
including the possibility $c=0$.

For each such choice, $S=B-cc^T$ is fixed. The set of $A$ solving
$A^TD_\lambda A=S$ is a coset of the orthogonal centralizer of
$D_\lambda$, which is
\[
 O(2)\times O(2)\times\{\pm1\}^{d-4}.
\]
Indeed, a commuting orthogonal map preserves each distinct eigenspace,
and any orthogonal map on each eigenspace commutes. Each $O(2)$
has dimension one, so this centralizer has dimension two. The
fiber is contained in a finite union of such cosets, proving the
claim.
\end{proof}

\subsection{Signed orbit strata}
Index positive roots by edges $e=(i,j)$ with $i<j$, ordered
lexicographically.
For $A\notin\Exc$, put two edges in the same class when
\[
 A\alpha_e=\sigma h_\theta A\alpha_{e'}
 \quad\text{for some }\sigma\in\{\pm1\},\ \theta\in\R.
\]
By \cref{lem:orbits}, the sign relative to a fixed representative is
unique. Choose the least edge in each class as its representative,
and give that representative sign $+1$. There are finitely many
resulting signed partitions.

Fix one signed partition, with $m$ classes and representatives
$e_1,\ldots,e_m$, and let $Z\subset\Orth\setminus\Exc$ be the set
of placements realizing it exactly. It is semialgebraic: relations
and nonrelations are expressed using equalities and inequalities
between values of \eqref{eq:orbit-map}. It is also preserved by
left multiplication by every $h_\theta$. The word \emph{stratum} here
means a piece of this finite semialgebraic partition; no smoothness,
connectedness, or constant orbit dimension for the individual roots is
assumed. All subsequent dimension arguments use \cref{lem:dimension},
not a smooth transversality theorem.
Define
\begin{align}
 L_Z&=\{c\in V:\ip{c}{\alpha_e}
       =\sigma_e\ip{c}{\alpha_{e_a}}
       \text{ for }e\text{ in class }a\},
       \label{eq:LZ}\\
 \mathscr S_Z&=\{B\in\Sym(V):
       \ip{\alpha_e}{B\alpha_e}
       =\ip{\alpha_{e_a}}{B\alpha_{e_a}}
       \text{ for }e\text{ in class }a\}.
       \label{eq:SZ}
\end{align}
Write $\ell=\dim L_Z$. By \cref{lem:quadratic},
\begin{equation}\label{eq:Sdim}
 \dim\mathscr S_Z=m.
\end{equation}

\begin{lemma}[Dimension of a signed stratum]\label{lem:stratum}
For every nonempty signed stratum $Z$,
\begin{equation}\label{eq:stratum-bound}
 \dim Z+\ell\leq m.
\end{equation}
\end{lemma}
\begin{proof}
For $A\in Z$, $c\in L_Z$, and the parameters in \eqref{eq:Dlambda},
we have
\begin{equation}\label{eq:image-in-S}
 A^TD_\lambda A+cc^T\in\mathscr S_Z.
\end{equation}
For the first term, this is because $D_\lambda$ commutes with
$h_\theta$ and quadratic evaluations do not see the sign $\sigma_e$.
For the second term, it follows by squaring the defining equalities
of $L_Z$.

Restrict \eqref{eq:spectral-map} to
\[
 Z\times(0,1)\times(4,5)
       \times\bigl(L_Z\cap\{\norm c<1\}\bigr).
\]
The domain has dimension $\dim Z+2+\ell$, and its image lies in
the $m$-dimensional space \eqref{eq:SZ}. Its fibers have dimension
at most two by \cref{lem:spectral}. Applying
\eqref{eq:fiber-upper} yields
$\dim Z+2+\ell\leq m+2$, which is \eqref{eq:stratum-bound}.
The restriction $\norm c<1$ is used only for this auxiliary
spectral argument; it places no bound on the eventual linear
functional in the agreement problem.
\end{proof}

\begin{proposition}[Generic nonexceptional avoidance]\label{prop:generic}
The set
\[
 \mathcal B=\{F\in\Pol:\mathcal R_F(A)=0
                      \text{ for some }A\notin\Exc\}
\]
is semialgebraic and has dimension at most $\dim\Pol-1$.
In particular, $\Pol\setminus\mathcal B$ is dense in $\Pol$.
\end{proposition}
\begin{proof}
Fix a nonempty stratum $Z$ as above, and put $p=\dim\Pol$.
For each $A\in Z$, the orbits of the representatives and their
negatives are all distinct. Thus \cref{lem:interpolation} makes
the $m$-value evaluation map surjective. Consider the semialgebraic
incidence set
\begin{equation}\label{eq:incidence}
 \mathcal I_Z=\{(A,c,F)\in Z\times L_Z\times\Pol:
 F(A\alpha_{e_a})=\ip{c}{\alpha_{e_a}}
                         \ (1\leq a\leq m)\}.
\end{equation}
Over every $(A,c)\in Z\times L_Z$, its fiber is a nonempty affine
space of dimension $p-m$. By \cref{lem:dimension,lem:stratum},
\begin{equation}\label{eq:incidence-dim}
 \dim\mathcal I_Z=\dim Z+\ell+p-m\leq p.
\end{equation}

Every nonempty fiber of the projection $\mathcal I_Z\to\Pol$
contains the circle
\[
 \{(h_\theta A,c,F):\theta\in\R/2\pi\mathbb Z\}.
\]
It lies in the same fiber because $F$ is invariant and $Z$ is
preserved by the action. It is one-dimensional: if
$h_\theta A=A$, the surjectivity of $A$ gives $h_\theta=I$, and
its action on $E_1$ gives $\theta=0$ modulo $2\pi$.
The circle is semialgebraic by the polynomial parametrization
preceding \cref{def:exceptional}. Consequently
\eqref{eq:fiber-lower} and \eqref{eq:incidence-dim} bound the
projection's dimension by $p-1$.

A polynomial has exact labels at a placement in $Z$ if and only
if it occurs in this projection. For the forward implication,
\cref{lem:exact} supplies $c$, and signed orbit invariance forces
$c\in L_Z$. For the converse, the representative equations in
\eqref{eq:incidence} extend to every edge by oddness and circle
invariance. Taking the finite union of these projections over
all signed partitions gives precisely $\mathcal B$. The final
assertion follows because a smaller-dimensional semialgebraic
set has empty interior.
\end{proof}

\section{The cubic identities}\label{sec:cubic}

Write $z_1=x+iy$, $z_2=u+iv$. Then
\begin{equation}\label{eq:real-seed}
 Q=(x^2-y^2)u+2xyv,\qquad
 H=(x^2+y^2)\bigl(2xyu-(x^2-y^2)v\bigr).
\end{equation}
Let
\[
 J(z_1,z_2,t)=(iz_1,2iz_2,0),\qquad
 X(A)_{ij}=\ip{p_i}{Jp_j}.
\]
For an odd function $G$ define
\begin{equation}\label{eq:Lambda}
 \Lambda_A(G)=\sum_{i<j}X(A)_{ij}
                     G\!\left(\frac{p_i-p_j}{\sqrt2}\right).
\end{equation}
The matrix $X(A)$ is skew-symmetric and $X(A)\one=0$ by
\eqref{eq:frame}. Hence \eqref{eq:cycle} gives
\begin{equation}\label{eq:exact-Lambda}
 \mathcal R_G(A)=0\quad\Longrightarrow\quad\Lambda_A(G)=0.
\end{equation}

The seed and the annihilator used in this section are those of
\cite{ChatGPT2026}. We give direct proofs, including the normalization
constants, because the later argument uses them in every dimension.

\begin{lemma}[Cubic expansion]\label{lem:cubic-expansion}
For any real homogeneous cubic $C$ and vectors $a,b$,
\begin{equation}\label{eq:cubic-expansion}
 C(a-b)=C(a)-C(b)-\ip{\nabla C(a)}b+\ip a{\nabla C(b)}.
\end{equation}
\end{lemma}
\begin{proof}
Write $C(w)=T(w,w,w)$ for its symmetric trilinear polarization.
Expansion of $T(a-b,a-b,a-b)$ gives mixed terms
$-3T(a,a,b)+3T(a,b,b)$, which are exactly the two gradient terms
in \eqref{eq:cubic-expansion}.
\end{proof}

\begin{lemma}[Cubic annihilation]\label{lem:annihilator}
For every $A\in\Orth$,
\begin{equation}\label{eq:cubic-annihilation}
 \Lambda_A(Q)=0.
\end{equation}
\end{lemma}
\begin{proof}
Apply \eqref{eq:cubic-expansion} to \eqref{eq:Lambda}. The terms
$Q(p_i)-Q(p_j)$ are exact and make no contribution. Also,
by \eqref{eq:frame} and the skew-symmetry of $J$,
\[
 \sum_jX(A)_{ij}p_j=-Jp_i,\qquad
 \sum_iX(A)_{ij}p_i=Jp_j.
\]
Since the product of two skew-symmetric edge labels is unchanged
when $i,j$ are exchanged, the sum over $i<j$ is half the sum
over all ordered pairs. Therefore
\begin{align*}
 \Lambda_A(Q)
 &=\frac1{4\sqrt2}\sum_{i,j}X(A)_{ij}
    \bigl(-\ip{\nabla Q(p_i)}{p_j}
                 +\ip{p_i}{\nabla Q(p_j)}\bigr)\\
 &=\frac1{2\sqrt2}\sum_i\ip{\nabla Q(p_i)}{Jp_i}=0.
\end{align*}
The last equality follows by differentiating
$Q(h_\theta w)=Q(w)$ at $\theta=0$.
\end{proof}

Let $x,y,u,v\in\R^N$ now denote the first four coordinate
\emph{columns} of the simplex frame $(p_i)$. All products of such
columns below are componentwise. Set
\begin{equation}\label{eq:P-and-G}
 P=(x\ \ y\ \ u\ \ v),\qquad
 G=(2xu+2yv\ \ -2yu+2xv\ \ x^2-y^2\ \ 2xy).
\end{equation}
Thus the rows of $G$ are the first four gradient coordinates of $Q$.

\begin{lemma}[Symmetric gradient matrix]\label{lem:gradient}
If $\mathcal R_Q(A)=0$, there is a symmetric real $4\times4$ matrix
$S$ such that
\begin{equation}\label{eq:gradient-matrix}
 G=PS.
\end{equation}
\end{lemma}
\begin{proof}
By \cref{lem:frame}, $P^TP=I_4$ and $\one^TP=0$.
The same identities give $\one^TG=0$: the mixed products have
zero sum, and $\sum_i(x_i^2-y_i^2)=1-1=0$.
Applying \eqref{eq:cubic-expansion} to every edge and projecting
away the exact term gives
\begin{equation}\label{eq:cubic-residual}
 \mathcal R_Q(A)=\frac1{2\sqrt2}(PG^T-GP^T).
\end{equation}
Exactness therefore gives $PG^T=GP^T$. Multiplication on the right
by $P$ yields $G=P(G^TP)$, while multiplication on both sides by
$P^T$ and $P$ gives $G^TP=P^TG$. Thus $S=G^TP$ is symmetric and
satisfies \eqref{eq:gradient-matrix}.
\end{proof}

\begin{lemma}[Allowed normalizations]\label{lem:normalizations}
Each of the following operations preserves exactness of the $Q$
labels and preserves the number $\Lambda_A(H)$:
\begin{enumerate}[label=(\roman*)]
 \item left multiplication of $A$ by $h_\theta$;
 \item simultaneous complex conjugation
       $C(z_1,z_2,t)=(\overline z_1,\overline z_2,t)$;
 \item a permutation of the simplex indices.
\end{enumerate}
\end{lemma}
\begin{proof}
The maps $h_\theta$ preserve $Q,H$ and commute with $J$.
For conjugation, $Q(Cw)=Q(w)$, $H(Cw)=-H(w)$, and $C^TJC=-J$.
Thus both factors of each summand of $\Lambda_A(H)$ change sign.
Finally, the exactness condition is invariant under relabeling.
The identity
\[
 \Lambda_A(H)=\frac12\sum_{i,j}X(A)_{ij}Y_H(A)_{ij}
\]
shows invariance under permutations as well. A permutation of
indices is realized by the corresponding orthogonal automorphism
of $V$, so all the resulting frames are still legitimate placements.
\end{proof}

\section{Exceptional cubic solutions and the quintic obstruction}
\label{sec:exceptional}

The argument in this section classifies only the information needed
at exceptional cubic solutions. It does \emph{not} classify every
solution of $\mathcal R_Q(A)=0$.

\begin{proposition}[Necessary normal form]\label{prop:normal-form}
Suppose $A\in\Exc$ and $\mathcal R_Q(A)=0$. After operations allowed
in \cref{lem:normalizations}, there are disjoint sets of indices
$K_1,K_2$, of even cardinalities $k_1,k_2\geq2$, and real numbers
$\gamma\ne0$ and $\omega$ such that
\begin{equation}\label{eq:normal-supports}
 x_i=\begin{cases}\pm k_1^{-1/2},&i\in K_1,\\0,&i\notin K_1,
       \end{cases}
 \qquad
 y_i=\begin{cases}\pm k_2^{-1/2},&i\in K_2,\\0,&i\notin K_2.
       \end{cases}
\end{equation}
On each support the plus and minus signs occur equally often, and
\begin{equation}\label{eq:normal-uv}
 u=\frac{x^2-y^2}{\gamma},\qquad
 v_i=\omega\quad(i\in K_1\cup K_2),\qquad
 \gamma^2=\frac1{k_1}+\frac1{k_2}.
\end{equation}
The complement of $K_1\cup K_2$ is nonempty. These are necessary
conditions only; the proposition does not claim that every frame
satisfying them is exceptional.
\end{proposition}
\begin{proof}
Choose a root in $E_1\cup E_2$ and distinguish the two possibilities.
Throughout the proof, the coordinate columns are orthonormal,
have zero mean, and satisfy \eqref{eq:gradient-matrix}.

\subsubsection*{Case 1: a root belongs to $E_1$}
A circle rotation and a relabeling make this root the real $z_1$
axis and its indices $0,1$. Since $A^Te_x=\alpha_{01}$, the
column $x$ is
\[
 x_0=\frac1{\sqrt2},\qquad x_1=-\frac1{\sqrt2},\qquad
 x_i=0\quad(i\ne0,1).
\]
The other coordinate columns have equal entries at indices $0,1$.
Write
\[
 y_0=y_1=\tau,\qquad u_0=u_1=U,\qquad v_0=v_1=W.
\]
Subtracting rows $0$ and $1$ in \eqref{eq:gradient-matrix} gives
\begin{equation}\label{eq:case1-first-row}
 (S_{xx},S_{xy},S_{xu},S_{xv})=(2U,2W,0,2\tau).
\end{equation}
Here the subscripts indicate the ordered coordinates $x,y,u,v$.
Since $2xy=2\tau x$, the last column of $G$ and the independence
of the columns of $P$ give
\[
 S_{xv}=2\tau,\qquad S_{yv}=S_{uv}=S_{vv}=0.
\]
Symmetry of $S$ and the remaining columns of
\eqref{eq:gradient-matrix} consequently give the following form of the symmetric coefficient matrix:
\[
 S=\begin{pmatrix}
  2U&2W&0&2\tau\\
  2W&-2a&-2b&0\\
  0&-2b&\gamma&0\\
  2\tau&0&0&0
 \end{pmatrix}.
\]
Here $xu=Ux$ and $xv=Wx$, because $x$ is supported at $0,1$.
The first three column equations therefore reduce to
\begin{align}
 yv&=Wy+\tau v,\label{eq:case1-yv}\\
 yu&=a y+b u,\label{eq:case1-yu}\\
 x^2-y^2&=-2b y+\gamma u,\label{eq:case1-squares}
\end{align}
for some real $a,b,\gamma$. Specifically,
$a=-S_{yy}/2$, $b=-S_{yu}/2$, and $\gamma=S_{uu}$.
Taking the inner product of \eqref{eq:case1-squares} with $u$
and using \eqref{eq:case1-yu} gives
\begin{equation}\label{eq:case1-gamma}
 \gamma=\ip{x^2}{u}-\ip{y^2}{u}
        =U-\ip{y}{yu}=U-a.
\end{equation}
This derives all coefficients used below.

We first prove $\tau=0$. At index $0$, \eqref{eq:case1-yv}
implies $\tau W=0$. If $\tau\ne0$, then $W=0$ and
$(y-\tau\one)v=0$. Thus $v$ is supported where $y_i=\tau$.
The unit column $v$ vanishes at $0,1$, so there exists an index
$i\notin\{0,1\}$ with $v_i\ne0$ and $y_i=\tau$.
Comparing \eqref{eq:case1-squares} at $0$ and $i$ gives
\[
 \gamma(U-u_i)=\frac12.
\]
In particular $U\ne u_i$. Comparison of \eqref{eq:case1-yu}
at these indices gives $(\tau-b)(U-u_i)=0$, so $b=\tau$.
Substitution at index $0$ then gives $a=0$.
Equation \eqref{eq:case1-yu} becomes $yu=\tau u$.
For an index $j\notin\{0,1\}$ with $y_j\ne\tau$, this forces
$u_j=0$, and \eqref{eq:case1-squares} gives
$-y_j^2=-2\tau y_j$. Therefore every entry of $y$ lies in
$\{0,\tau,2\tau\}$. Since $y_0=\tau\ne0$, all nonzero entries
have the same sign, contradicting $\sum_jy_j=0$.
This proves $\tau=0$.

At index $0$, \eqref{eq:case1-squares} now gives $\gamma U=1/2$,
so $\gamma,U\ne0$. Equation \eqref{eq:case1-yu} at that index
gives $b=0$. Hence
\begin{equation}\label{eq:case1-reduced}
 u=\frac{x^2-y^2}{\gamma},\qquad yu=a y.
\end{equation}
The supports of $x$ and $y$ are disjoint because $y_0=y_1=0$.
On the nonempty support of $y$, the second equation says $u_i=a$;
the first says $y_i^2=-\gamma a$. This is a fixed positive value.
Because $y$ has zero mean and norm one, its support has even
cardinality and its nonzero entries have equal positive and negative
multiplicities, with magnitude the reciprocal square root of that
cardinality. Equation \eqref{eq:case1-yv} now says $v_i=W$ on
this support; the same holds on $\{0,1\}$ by definition.
This is \eqref{eq:normal-supports}--\eqref{eq:normal-uv}, with
$K_1=\{0,1\}$ and $\omega=W$, apart from the final norm identity,
which is proved after Case 2.

\subsubsection*{Case 2: a root belongs to $E_2$}
A circle rotation and relabeling make the root the real $z_2$
axis, with indices $0,1$. Thus
\[
 u_0=\frac1{\sqrt2},\qquad u_1=-\frac1{\sqrt2},\qquad
 u_i=0\quad(i\ne0,1).
\]
Write $x_0=x_1=X$ and $y_0=y_1=Y$.
Subtracting rows in \eqref{eq:gradient-matrix} gives the $u$ row
of $S$ as $(2X,-2Y,0,0)$. Symmetry therefore gives
\[
 x^2-y^2=2Xx-2Yy.
\]
At index $0$, this implies $X^2=Y^2$. If necessary, simultaneous
complex conjugation changes $Y$ to $-Y$, so assume $X=Y$.
Put
\begin{equation}\label{eq:rs-def}
 r=\frac{x-y}{\sqrt2},\qquad
 s=\frac{x+y}{\sqrt2},\qquad \eta=\sqrt2X.
\end{equation}
These are orthonormal columns and satisfy
\begin{equation}\label{eq:case2-basic}
 rs=\eta r,\qquad ru=0,\qquad su=\eta u,
 \qquad r_0=r_1=0,\quad s_0=s_1=\eta.
\end{equation}

In the orthogonal coordinates $r,s,u,v$ the cubic becomes
\[
 Q=2rsu+(s^2-r^2)v,
\]
and its gradient columns become
\[
 (2su-2rv,\quad 2ru+2sv,\quad 2rs,\quad s^2-r^2).
\]
This is a passive change of coordinates: the displayed polynomial is
$Q$ expressed in $r,s,u,v$, not an assertion that $Q$ or $H$ is invariant
under this change. If $T$ is the corresponding orthogonal coordinate
matrix, the frame and gradient matrices become $PT$ and $GT$, and
$G=PS$ becomes $GT=(PT)(T^TST)$. Thus the coefficient matrix remains
symmetric. The actual normalization of the placement at the end of this
case will use the allowed circle rotation $h_{\pi/4}$. By \eqref{eq:case2-basic}, its third gradient column
is $2\eta r$. Symmetry fixes the $u$ coefficients in the other
columns. The additional identities
\[
 \ip{rv}{s}=\ip{sv}{r}=\sum_i r_is_iv_i
                         =\eta\ip r v=0
\]
remove their cross coefficients. It follows that there are real
$a,c,e,f,\gamma$ with
\begin{align}
 rv&=a r+c v,& sv&=e s+f v,\label{eq:case2-products}\\
 s^2-r^2&=-2c r+2f s+\gamma v.
                  \label{eq:case2-squares}
\end{align}
Equivalently, in the ordered coordinates $r,s,u,v$, the coefficient
matrix has the form
\[
 \begin{pmatrix}
 -2a&0&2\eta&-2c\\
 0&2e&0&2f\\
 2\eta&0&0&0\\
 -2c&2f&0&\gamma
 \end{pmatrix}.
\]
This displays explicitly the coefficients $-2c$ and $2f$ forced by
symmetry. Taking the inner product of
\eqref{eq:case2-squares} with $v$ gives
\begin{equation}\label{eq:case2-gamma}
 \gamma=\ip s{sv}-\ip r{rv}=e-a.
\end{equation}

We claim $\eta=0$. Suppose otherwise, and let $K=\supp r$.
It is nonempty, and $s_i=\eta$ for $i\in K$.
If $c\ne0$, the first equation of \eqref{eq:case2-products}
forces $v_i=0$ off $K$. Thus $sv=\eta v$, so independence of
$s,v$ in the second equation gives $e=0$ and $f=\eta$.
Off $K$, \eqref{eq:case2-squares} now gives
$s_i^2=2\eta s_i$. On $K$, one has $s_i=\eta$.
All entries of $s$ therefore belong to $\{0,\eta,2\eta\}$,
contradicting its zero mean. Hence $c=0$.

The first equation of \eqref{eq:case2-products} now gives $v_i=a$
on $K$. Let $v_0=v_1=W$. Comparing
\eqref{eq:case2-squares} at $0$ and an index $i\in K$ yields
\[
 r_i^2=\gamma(W-a)>0.
\]
In particular $W\ne a$. Comparing the second equation of
\eqref{eq:case2-products} at those same indices gives
$(\eta-f)(W-a)=0$. Therefore $f=\eta$, and substitution gives
$e=0$. Now $(s-\eta\one)v=0$. At any index outside $K$ where
$s_i\ne\eta$, one has $v_i=0$, and
\eqref{eq:case2-squares} yields $s_i\in\{0,2\eta\}$.
All other entries equal $\eta$. This gives the same zero-mean
contradiction. Thus $\eta=0$.

We now have $rs=0$. Multiplying the first equation of
\eqref{eq:case2-products} by $s$ gives $ce=cf=0$;
multiplying the second by $r$ gives $fa=fc=0$, since the relevant
columns are linearly independent. If $c\ne0$, then $e=f=0$, so
$sv=0$. At an index where $s_i\ne0$, both $r_i$ and $v_i$ vanish,
and \eqref{eq:case2-squares} would give $s_i^2=0$, a contradiction.
Hence $c=0$. Similarly, if $f\ne0$, then $a=c=0$ and $rv=0$;
at an index where $r_i\ne0$, the same equation gives
$-r_i^2=0$. Therefore $f=0$ as well. We have proved
\begin{equation}\label{eq:case2-reduced}
 rv=a r,\qquad sv=e s,\qquad s^2-r^2=\gamma v.
\end{equation}
The number $\gamma$ is nonzero: otherwise $r^2=s^2$, impossible
for nonzero columns with disjoint supports.

On the support of $r$ one has $v_i=a$ and $r_i^2=-\gamma a$;
on the support of $s$ one has $v_i=e$ and $s_i^2=\gamma e$.
Thus each nonzero magnitude is constant on its support. Zero mean
and unit norm give even support sizes and balanced signs as in
\eqref{eq:normal-supports}. Outside the two supports, $v_i=0$.

Finally apply the allowed circle rotation $h_{\pi/4}$. By
\eqref{eq:rs-def}, the new complex columns are
\[
 z_1'=r+is,\qquad z_2'=-v+iu.
\]
Thus the new real $z_2$ column is $(r^2-s^2)/\gamma$.
The new imaginary column is the old $u$, which is supported at
indices $0,1$, where $r=s=0$. Consequently it is zero on both
supports. This proves the asserted form with $\omega=0$.

\subsubsection*{The final norm identity}
In either case, disjointness of the supports and the fact that
$u$ has norm one give
\[
 1=\norm u^2
  =\frac1{\gamma^2}
     \left(\sum_i x_i^4+\sum_i y_i^4\right)
  =\frac1{\gamma^2}\left(\frac1{k_1}+\frac1{k_2}\right).
\]
Finally, the complement $O$ of $K_1\cup K_2$ is nonempty. Otherwise
$v=\omega\one$ would have zero mean only if $v=0$, contradicting its
unit norm. In particular $k_1+k_2\leq N-1$.
This completes the proof of \eqref{eq:normal-uv} and the proposition.
\end{proof}

\begin{lemma}[Evaluation in the normal form]\label{lem:normal-evaluation}
For a frame in the form of \cref{prop:normal-form},
\begin{equation}\label{eq:nonzero-obstruction}
 \Lambda_A(H)=\frac1{\sqrt2\,k_1k_2\gamma}\ne0.
\end{equation}
\end{lemma}
\begin{proof}
For an oriented edge $(i,j)$ write
$\Delta x=x_i-x_j$ and similarly for the other coordinates.
The normalization by $\sqrt2$ gives
\begin{equation}\label{eq:H-edge}
 H\!\left(\frac{p_i-p_j}{\sqrt2}\right)
 =\frac{\bigl((\Delta x)^2+(\Delta y)^2\bigr)
     \bigl(2\Delta x\Delta y\Delta u
        -((\Delta x)^2-(\Delta y)^2)\Delta v\bigr)}{4\sqrt2}.
\end{equation}
The generator gives
\begin{equation}\label{eq:X-edge}
 X(A)_{ij}=-x_i y_j+y_i x_j-2u_i v_j+2v_i u_j.
\end{equation}
Let $O$ be the complement of $K_1\cup K_2$, and put
\[
 U_1=\frac1{k_1\gamma},\qquad
 U_2=-\frac1{k_2\gamma}.
\]
Then $u_i=U_a$ on $K_a$, $u_i=0$ on $O$, and
$U_1-U_2=\gamma$ by \eqref{eq:normal-uv}.
We group the edge contributions. The product of the two factors
in \eqref{eq:Lambda} is independent of edge orientation, so each
edge may be oriented to agree with the indicated grouping.

\subsubsection*{Edges between $K_1$ and $K_2$}
For $i\in K_1$, $j\in K_2$, put $a=x_i$ and $b=y_j$.
Equations \eqref{eq:H-edge} and \eqref{eq:X-edge} give
\[
 H\!\left(\frac{p_i-p_j}{\sqrt2}\right)
    =-\frac{\gamma^3ab}{2\sqrt2},\qquad
 X(A)_{ij}=-ab-2\omega\gamma.
\]
Consequently the sum over these edges is
\begin{align}
 \frac{\gamma^3}{2\sqrt2}
 \sum_{i\in K_1,j\in K_2}
       \bigl(x_i^2y_j^2+2\omega\gamma x_i y_j\bigr)
 &=\frac{\gamma^3}{2\sqrt2}.
 \label{eq:cross-sum}
\end{align}
We used $\sum_{K_1}x_i^2=\sum_{K_2}y_j^2=1$ and
$\sum_{K_1}x_i=\sum_{K_2}y_j=0$.

\subsubsection*{Edges from $K_1$ to $O$}
Fix $j\in O$ and write $w=v_j$. For $i\in K_1$,
\[
 H\!\left(\frac{p_i-p_j}{\sqrt2}\right)
       =-\frac{\omega-w}{4\sqrt2\,k_1^2},\qquad
 X(A)_{ij}=-2U_1w.
\]
Summing over the $k_1$ possible indices $i$ gives
\begin{equation}\label{eq:K1O}
 \frac{w(\omega-w)}{2\sqrt2\,k_1^2\gamma}.
\end{equation}

\subsubsection*{Edges from $K_2$ to $O$}
For the same fixed $j$ and $i\in K_2$,
\[
 H\!\left(\frac{p_i-p_j}{\sqrt2}\right)
       =\frac{\omega-w}{4\sqrt2\,k_2^2},\qquad
 X(A)_{ij}=-2U_2w.
\]
Their sum is
\begin{equation}\label{eq:K2O}
 \frac{w(\omega-w)}{2\sqrt2\,k_2^2\gamma}.
\end{equation}

\subsubsection*{All remaining edges and the total}
Edges internal to $K_1$ or to $K_2$ have $\Delta u=\Delta v=0$,
so their $H$ labels vanish. Edges internal to $O$ have
$\Delta x=\Delta y=0$, so their labels also vanish.
Because $v$ has zero mean and norm one,
\[
 \sum_{j\in O}v_j=-(k_1+k_2)\omega,\qquad
 \sum_{j\in O}v_j^2=1-(k_1+k_2)\omega^2.
\]
It follows that
\begin{equation}\label{eq:outside-moment}
 \sum_{j\in O}v_j(\omega-v_j)=-1.
\end{equation}
Thus \eqref{eq:cross-sum}--\eqref{eq:outside-moment} give
\begin{align*}
 \Lambda_A(H)
 &=\frac1{2\sqrt2}
       \left(\gamma^3-
          \frac{k_1^{-2}+k_2^{-2}}{\gamma}\right)\\
 &=\frac{(k_1^{-1}+k_2^{-1})^2-k_1^{-2}-k_2^{-2}}
          {2\sqrt2\,\gamma}
  =\frac1{\sqrt2\,k_1k_2\gamma}.
\end{align*}
The denominator is nonzero. This proves the assertion.
\end{proof}

\begin{corollary}[Exceptional obstruction]\label{cor:exceptional}
If $A\in\Exc$ and $\mathcal R_Q(A)=0$, then $\Lambda_A(H)\ne0$.
\end{corollary}
\begin{proof}
Apply \cref{prop:normal-form,lem:normal-evaluation}, and use the
invariance of $\Lambda_A(H)$ under the normalizations in
\cref{lem:normalizations}.
\end{proof}

\section{Completion of the counterexample construction}\label{sec:completion}

\begin{proposition}[An open set avoiding exceptional placements]
\label{prop:exceptional-open}
There exists $\eps_0>0$ such that, for every
$0<|\eps|<\eps_0$, the polynomial $F_\eps=Q+\eps H$ satisfies
\[
 \mathcal R_{F_\eps}(A)\ne0\qquad(A\in\Exc).
\]
For each such $\eps$, this property persists in a neighborhood of
$F_\eps$ in $\Pol$.
\end{proposition}
\begin{proof}
If the first assertion failed, there would exist nonzero
$\eps_j\to0$ and $A_j\in\Exc$ such that
$\mathcal R_{Q+\eps_jH}(A_j)=0$.
By \eqref{eq:exact-Lambda} and \cref{lem:annihilator},
\[
 0=\Lambda_{A_j}(Q+\eps_jH)
   =\eps_j\Lambda_{A_j}(H).
\]
Hence $\Lambda_{A_j}(H)=0$ for every $j$.
Compactness gives a subsequence converging to
$A_\infty\in\Exc$. The residual is continuous, so
$\mathcal R_Q(A_\infty)=0$; likewise,
$\Lambda_{A_\infty}(H)=0$. This contradicts
\cref{cor:exceptional}.

Fix an admissible $\eps$. Continuity and compactness give
\begin{equation}\label{eq:positive-gap}
 \delta=\min_{A\in\Exc}
             \norm{\mathcal R_{F_\eps}(A)}_{\mathrm F}>0,
\end{equation}
where $\norm{\cdot}_{\mathrm F}$ is the Frobenius matrix norm.
Because $\Pol$ is finite-dimensional and $\Orth$ is compact,
there is a constant $C>0$ such that
\begin{equation}\label{eq:uniform-bound}
 \sup_{A\in\Orth}\norm{\mathcal R_R(A)}_{\mathrm F}
                                  \leq C\norm R\qquad(R\in\Pol).
\end{equation}
For example, expand $R$ in a fixed basis of $\Pol$, bound each
basis residual on $\Orth$, and use equivalence of norms in a
finite-dimensional vector space. By linearity of the residual,
every $F_\eps+R$ with $\norm R<\delta/(2C)$ has residual norm
at least $\delta/2$ on $\Exc$. This proves the openness assertion.
\end{proof}

\begin{proof}[Proof of \Cref{thm:perturbative,thm:main}]
Fix $0<|\eps|<\eps_0$ as in \cref{prop:exceptional-open}, and
let $\eta>0$. By \cref{prop:generic}, the complement of
$\mathcal B$ is dense in $\Pol$. The open ball about $F_\eps$
of radius $\min\{\eta,\delta/(2C)\}$ therefore contains a
polynomial $F\notin\mathcal B$. Write $R=F-F_\eps$.
Then $\norm R<\eta$, and \cref{prop:exceptional-open} excludes
exact labels on $\Exc$. The definition of $\mathcal B$ excludes
exact labels everywhere else. Thus
$\mathcal R_F(A)\ne0$ for every $A\in\Orth$.
By \cref{lem:exact}, this is precisely the failure of
\eqref{eq:agreement}. The polynomial $F\in\Pol$ is odd and
has degree at most $D$, and its sphere restriction is smooth.
This proves both theorems.
\end{proof}

\begin{remark}\label{rem:scope}
The argument is an existence proof in an explicitly specified
finite-dimensional polynomial space. It does not exhibit coefficients
for $R$, a numerical $\eps_0$, or an optimal degree bound. It also
does not prove that $Q+\eps H$ alone is a counterexample for every
$d\geq4$. The purpose of $R$ is to avoid the nonexceptional
placements without classifying all their cubic solutions.
\end{remark}

\section{The constant-width and universal-cover consequence}
\label{sec:geometry}

The passage from odd functions to small perturbations of the support
function of a ball is standard; see
\cite[Section 1, Proposition 1]{Kuperberg1999} and
\cite[Lemma 8]{Chen2026}. We include the needed implication with
width normalized to one.

\begin{lemma}[Realizing a small odd perturbation]\label{lem:support}
Let $f:S^{d-1}\to\R$ be smooth and odd. For all sufficiently small
$\tau>0$, there is a compact convex body $K$ of constant width one
whose support function on the sphere is
\begin{equation}\label{eq:support}
 h_K(w)=\frac12+\tau f(w).
\end{equation}
\end{lemma}
\begin{proof}
For $w\ne0$ define
\[
 g(w)=\norm w\,f\!\left(\frac w{\norm w}\right),\qquad
 \widehat h(w)=\tfrac12\norm w+\tau g(w),
\]
and set $\widehat h(0)=0$. Both functions are homogeneous of degree
one for positive scalars. Differentiating Euler's identity gives
$\nabla^2g(w)w=0$ and $\nabla^2\widehat h(w)w=0$.
At $\norm w=1$,
\[
 \nabla^2\widehat h(w)
    =\tfrac12(I-ww^T)+\tau\nabla^2g(w).
\]
The Hessian of $g$ is uniformly bounded on the unit sphere. Thus,
for sufficiently small $\tau>0$, the Hessian of $\widehat h$
is positive definite on $w^\perp$ and positive semidefinite on
$\R^d$ at each nonzero $w$. Here radial and mixed terms vanish,
and homogeneity extends the assertion from the unit sphere to
all $w\ne0$.

The restriction of $\widehat h$ to a line not passing through the
origin is convex by the nonnegative second derivative. On a line
through the origin, it is piecewise linear, and the right slope
minus the left slope, using a unit direction $w$, is
\[
 \widehat h(w)+\widehat h(-w)=1.
\]
Hence it is convex on those lines as well. It follows that
$\widehat h$ is convex on all of $\R^d$.

Choose $\tau$ still smaller if needed so that
$\tau\norm f_\infty<1/2$. Define
\[
 K=\{z\in\R^d:\ip z w\leq\widehat h(w)
                          \text{ for every }w\in S^{d-1}\}.
\]
This is closed and convex, contains a ball of positive radius,
and is bounded by $\max_{S^{d-1}}\widehat h$.
To verify that its support function is exactly $\widehat h$,
fix $w\ne0$. Convexity and differentiability at $w$ imply
\[
 \widehat h(y)\geq\widehat h(w)
     +\ip{\nabla\widehat h(w)}{y-w}
     =\ip{\nabla\widehat h(w)}y,
\]
where the last equality is Euler's identity.
Thus $\nabla\widehat h(w)\in K$, and
$\ip{\nabla\widehat h(w)}w=\widehat h(w)$.
The defining inequality of $K$ gives the reverse support bound.
Therefore $h_K=\widehat h$ and \eqref{eq:support} holds.
Finally, oddness gives $h_K(w)+h_K(-w)=1$ for every unit direction,
which is constant width one.
\end{proof}

\begin{proof}[Proof of \Cref{cor:cover}]
Take $f=F|_{S^{d-1}}$ from \cref{thm:main}, and form $K$ as in
\cref{lem:support}. Suppose $K\subset A C_d+b$.
For each $\alpha=\alpha_{ij}$, the two defining half-spaces give
\[
 h_K(A\alpha)\leq\tfrac12+\ip b{A\alpha},\qquad
 h_K(-A\alpha)\leq\tfrac12-\ip b{A\alpha}.
\]
The two left sides sum to one, as do the right sides. Both
inequalities must therefore be equalities, so
\[
 \tau f(A\alpha)=\ip b{A\alpha}
 \quad\text{for every positive root }\alpha.
\]
Division by $\tau>0$ contradicts \cref{thm:main}.

For completeness, a compact convex body of constant width one has
diameter one. For distinct $x,y\in K$, use
$w=(x-y)/\norm{x-y}$ to get
$\norm{x-y}\leq h_K(w)+h_K(-w)=1$.
Conversely, support points in opposite directions have scalar
separation one, so their distance is at least one. Thus this $K$
is itself a set of diameter one not covered by any congruent copy
of $C_d$.
\end{proof}

\appendix
\section{Verification notes and supplementary algebra checks}
\label{app:verification}

The central dependencies of the proof can be checked independently.
\Cref{lem:spectral} is a statement about real symmetric matrices;
it uses neither the cubic nor the exceptional normal form.
\Cref{lem:stratum} applies it to the quadratic evaluation
isomorphism. \Cref{prop:generic} then uses only the interpolation
lemma and semialgebraic fiber dimensions. On the other side,
\cref{prop:normal-form} is a consequence of the four-column
orthonormality relations and the symmetric gradient matrix.
The explicit sum in \cref{lem:normal-evaluation} finishes the
exceptional analysis. Compactness combines the two parts.

For checking \cref{lem:spectral}, the strict condition $\norm c<1$
and the separation of the eigenvalue clusters are essential: they
ensure that any multiple eigenvalue of $B$ is an unchanged double
eigenvalue of $S$, where the residue in \eqref{eq:secular} vanishes.
In the incidence argument, the potential $c$ is \emph{not} restricted
to that unit ball. The bounded auxiliary parameter is used only
to obtain \eqref{eq:stratum-bound}.

For checking the normal form, the coefficients
$\gamma=U-a$ in Case 1 and $\gamma=e-a$ in Case 2 follow from
inner products, not just from subtracting the two special rows.
The transformations in Case 2 include simultaneous conjugation,
which reverses both $H$ and the annihilating matrix and therefore
preserves their pairing. The factor $(\sqrt2)^{-5}=1/(4\sqrt2)$
in \eqref{eq:H-edge} is retained throughout.

The companion script \texttt{verify\_algebra\_revised.py} checks, using exact
symbolic arithmetic, the infinitesimal circle invariance of $Q,H$,
the cubic expansion, the normalized quintic edge formula, the
three nonzero edge-group sums, and their final simplification.
It also checks several explicitly specified four-column frames
satisfying the normal-form hypotheses, including zero and nonzero
$\omega$ and nonconstant outside entries of $v$. Additional checks cover
the quadratic-evaluation inverse, the general cubic-residual identity
on deterministic simplex frames, and characteristic-polynomial
identities for selected rank-one spectral examples. These are regression
checks for formulas: neither finitely many test frames nor the script
certifies exhaustiveness of the normal form or any dimension
argument. No computational output is used as a premise of the
proof.

\medskip
\noindent\textbf{Revision note (September 4, 2026).}
This revision corrects the scope of the unit-vector normalization statement,
makes the polar scale and universal-cover convention explicit, expands the
spectral multiplicity and normal-form coefficient arguments, and clarifies
the passive coordinate change in Case 2. The theorem statements and the
existential nature of the perturbation are unchanged. The accompanying
review and additional symbolic checks are AI-assisted and do not constitute
independent human review or a formal proof-assistant verification.

\begin{thebibliography}{JLYY11}

\bibitem[BCR98]{BCR1998}
Jacek Bochnak, Michel Coste, and Marie-Fran\c{c}oise Roy,
\emph{Real Algebraic Geometry}, Ergebnisse der Mathematik und ihrer
Grenzgebiete, 3. Folge, vol.~36, Springer-Verlag, Berlin, 1998.
\href{https://doi.org/10.1007/978-3-662-03718-8}
{doi:10.1007/978-3-662-03718-8}.

\bibitem[CG26]{ChatGPT2026}
ChatGPT,
\emph{Counterexamples to Makeev's regular-simplex conjecture in
dimensions four and five}, AI-generated, unrefereed repository
manuscript, August 2026, in the repository
\texttt{simplex4/math-results}, file
\texttt{makeev-code/a4/makeev\_d4\_d5.tex}.
\url{https://github.com/simplex4/math-results/blob/main/makeev-code/a4/makeev_d4_d5.tex}.
Consulted September 4, 2026. Authorship is reproduced as stated in
the manuscript; the repository account is not identified as a human
author here.

\bibitem[Ch26]{Chen2026}
Yizhen Chen,
\emph{Note on Lebesgue's universal cover problem},
arXiv:2607.27227v2 [math.MG], August 6, 2026.
\url{https://arxiv.org/abs/2607.27227v2}.

\bibitem[H80]{Hardt1980}
Robert M. Hardt,
\emph{Semi-algebraic local-triviality in semi-algebraic mappings},
American Journal of Mathematics \textbf{102} (1980), no.~2, 291--302.
\href{https://doi.org/10.2307/2374240}{doi:10.2307/2374240}.

\bibitem[HLTV11]{HLTV2011}
Robert Hardt, Pascal Lambrechts, Victor Turchin, and Ismar Voli\'c,
\emph{Real homotopy theory of semi-algebraic sets},
Algebraic \& Geometric Topology \textbf{11} (2011), no.~5, 2477--2545.
\href{https://doi.org/10.2140/agt.2011.11.2477}
{doi:10.2140/agt.2011.11.2477}.
Preprint: \url{https://arxiv.org/abs/0806.0476v3}.

\bibitem[JLYY11]{JiangLimYaoYe2011}
Xiaoye Jiang, Lek-Heng Lim, Yuan Yao, and Yinyu Ye,
\emph{Statistical ranking and combinatorial Hodge theory},
Mathematical Programming \textbf{127} (2011), no.~1, 203--244.
\href{https://doi.org/10.1007/s10107-010-0419-x}
{doi:10.1007/s10107-010-0419-x}.

\bibitem[K99]{Kuperberg1999}
Greg Kuperberg,
\emph{Circumscribing constant-width bodies with polytopes},
New York Journal of Mathematics \textbf{5} (1999), 91--100.
\url{https://nyjm.albany.edu/j/1999/5-6p.pdf}.
Preprint: \url{https://arxiv.org/abs/math/9809165v3}.

\bibitem[VM26]{VibeMathed2026}
VibeMathed,
\emph{Makeev's conjecture on universal cover}, online problem entry,
2026. \url{https://vibemathed.com/problem/makeev-s-conjecture-on-universal-cover}.
Consulted September 4, 2026. This entry links the repository
manuscript \cite{ChatGPT2026}; it is not used as a proof source.

\end{thebibliography}
\end{document}
